% SPDX-FileCopyrightText: 2026 Bart van de Lint
% SPDX-License-Identifier: AGPL-3.0-or-later
\section{How far the model can be trusted}
\label{app:trust}

This appendix holds the evidence behind the accuracy the main text states: how far the router converges and what forecasts cost it, the checks of the model against hand calculations and measurements, how much each part of the model moves the result, and the assumptions of the discussion that rest on data from outside the model.

\subsection{Convergence of the routing search}
\label{app:routing_convergence}

The cases are the 12 departures of the results for the \SI{140}{m} e-methanol ship and the Kamsarmax on VLSFO, each bare and with each rig, on the routes to Guadeloupe, New York and Tubarão.

\begin{table}[tbp]
\centering
\footnotesize
\caption{Largest change of the cost saving, in points, against the base settings, for the kite, the wing sails and the rotors, with the other settings at the base, at a corridor of stages \SI{360}{nmi} and points \SI{240}{nmi} apart. The last two rows, one figure each over all rigs, were measured on a finer search with knots \SI{48}{h} apart.}
\label{tab:routing_convergence}
\begin{tabular}{@{}p{10.5cm}r@{}}
\toprule
Setting changed & Change \\
\midrule
Margins 12\,\%, 20\,\% (kite 25\,\%, 50\,\%), bins of \SI{1.5}{h}, every second speed, 3 nodes & 0.72, 0.07, 0.48 \\
Time bins of \SI{1.5}{h} rather than \SI{3}{h} & 0.70, 0.05, 0.10 \\
Every second speed rather than every fourth & 0.29, 0.14, 0.22 \\
Edges of up to 3 lateral nodes rather than 2 & 0.74, 0.89, 1.12 \\
Margins 6\,\%, then 12\,\%, for ships without a kite & 0, 0.08, 0.01 \\
Weather every \SI{6}{h} rather than \SI{3}{h} & 0.56, 0.15, 0.45 \\
Table of 97 speeds, every fourth, against all 49 & 0.11 \\
Final pass in bins of \SI{0.125}{h}, or within \SI{48}{h} of the path & 0.22, 0.19 \\
\bottomrule
\end{tabular}
\end{table}

\Cref{tab:routing_grid} sets the error of the corridor spacing against the error of the search settings and against computing time, for the same cases, to the stages and points of the results (\SI{90}{nmi}, \SI{60}{nmi}) and the search settings of \cref{tab:numerics}, except that states are skipped against an estimate on knots \SI{48}{h} rather than \SI{24}{h} apart and, with a kite, within 10\,\% and then 20\,\% rather than 20\,\% and then 40\,\%. The settings of \cref{tab:numerics} take about 1.3 times the computing time of these and lose no more than 0.01 points of cost saving by skipping states, against up to 0.27 points with these. The reference is finer in both: stages \SI{45}{nmi} and points \SI{30}{nmi} apart, searched with bins of \SI{1.5}{h}, every second speed and the margins of the first row of \cref{tab:routing_convergence}, 13.6 times as slow as the results' corridor and search. Refining only the stages or only the points gains the kite little and lowers the wing sails' and rotors' saving, so the two are halved together. Halving the points alone also changes the headings a step can sail, since an edge reaches at most two points to either side: about 0, 34 and \SI{53}{\degree} off the corridor's line on the corridor of the results. Halving the points and letting an edge reach four of them keeps those headings and adds about 18 and \SI{45}{\degree} between them. On the corridor of the results this raises the kite's cost saving by 0.24 points on average and lowers the wing sails' by 0.10 and the rotors' by 0.01, by at most 0.6 points on one pair of ship and route, at 3.0 times the computing time; on stages \SI{180}{nmi} apart it moves them by 0.32, $-0.05$ and $-0.09$ points at 2.8 times. It lowers the cost of the ship without a rig by 1.5\,\% and 2.0\,\%, about as much as with a rig, so the headings between are not added. At the coarser corridors the search settings matter little against the spacing: at stages \SI{180}{nmi} apart the finer search recovers 0.2 of the kite's 2.2 points. At the corridor of the results the two errors are of a size: the finer search raises the kite's cost saving by 0.46 points and the rotors' by 0.11, half the spacing by 0.53 and 0.25, and both together by 1.26 and 0.48, more than the sum; none is corrected. The reference is itself not converged: with the finer search the kite's gain per halving of the spacing falls from 2.35 to 1.22 to 0.81 points and the rotors' from 1.75 to 0.71 to 0.37, so the kite's error falls more slowly than the rotors'.

\begin{table}[tbp]
\centering
\footnotesize
\caption{Change of the cost saving per tonne-kilometre, in points, against the reference, over the two ships and three routes: the mean over the six pairs of ship and route, and in brackets the largest change on one pair, for the kite, the wing sails and the rotors, with the computing time as a multiple of that of the base run of the results. Stages 360 and points \SI{120}{nmi} apart found no passage for one of the rotor ships, and stages 720 and points \SI{480}{nmi} apart none for the kite and rotor ships on one pair. The coarsest and the finest rows were timed on a later commit than the rest, which routes the same passages, each against the base run of its own commit.}
\label{tab:routing_grid}
\begin{tabular}{@{}llrccc@{}}
\toprule
Stages, points [nmi] & Search & Time & Kite & Wing sails & Rotors \\
\midrule
720, 480 & results & 0.13 & $-5.62$ (11.06) & $-1.81$ (5.20) & $-3.80$ (6.28) \\
360, 240 & results & 0.23 & $-4.64$ (8.21) & $-1.10$ (3.35) & $-2.86$ (5.70) \\
180, 120 & results & 0.46 & $-2.22$ (4.55) & $-0.58$ (2.50) & $-1.19$ (3.12) \\
90, 60 & results & 1.00 & $-1.26$ (2.39) & $-0.09$ (0.74) & $-0.48$ (1.72) \\
45, 30 & results & 1.88 & $-0.73$ (1.59) & $-0.14$ (0.31) & $-0.23$ (0.40) \\
180, 240 & results & 0.24 & $-3.70$ (5.79) & $-1.21$ (2.98) & $-2.72$ (3.91) \\
360, 120 & results & 0.36 & $-4.64$ (9.19) & $-1.68$ (4.04) & $-2.82$ (4.63) \\
90, 120 & results & 0.46 & $-2.15$ (4.06) & $-0.79$ (1.54) & $-1.58$ (2.30) \\
180, 60 & results & 0.52 & $-2.34$ (4.00) & $-1.10$ (2.94) & $-2.06$ (4.09) \\
\addlinespace
90, 60 & bins of \SI{1.5}{h} & 1.64 & $-0.94$ (1.96) & $-0.10$ (0.74) & $-0.34$ (1.59) \\
90, 60 & every second speed & 1.36 & $-1.18$ (2.43) & $-0.10$ (0.64) & $-0.55$ (1.40) \\
90, 60 & wider margins & 1.90 & $-1.14$ (2.32) & $-0.04$ (0.75) & $-0.41$ (1.41) \\
90, 60 & all three & 4.59 & $-0.81$ (1.77) & $-0.03$ (0.60) & $-0.37$ (1.33) \\
180, 120 & bins of \SI{1.5}{h} & 0.75 & $-2.12$ (4.61) & $-0.56$ (2.62) & $-1.20$ (3.19) \\
180, 120 & every second speed & 0.63 & $-2.29$ (5.30) & $-0.59$ (2.62) & $-1.23$ (3.06) \\
180, 120 & wider margins & 0.99 & $-2.16$ (4.55) & $-0.56$ (2.49) & $-1.13$ (3.12) \\
180, 120 & all three & 1.91 & $-2.03$ (4.64) & $-0.49$ (2.48) & $-1.07$ (3.03) \\
360, 240 & all three & 0.83 & $-4.38$ (7.79) & $-1.14$ (3.54) & $-2.82$ (5.57) \\
45, 30 & all three: the reference & 13.56 & $0.00$ (0.00) & $0.00$ (0.00) & $0.00$ (0.00) \\
\bottomrule
\end{tabular}
\end{table}

A table of 25 speeds, or 9 commanded, understates the kite's cost saving by 0.26 and 0.34 points on average and by up to 1.3 points. Whole knots from \SI{4}{kn} understate the kite's cost saving by 0.39 points on average, the rotors' by 0.15 and the wings' by 0.06.

On a corridor four times as coarse, with states skipped as in the results, the search settings lower the cost saving per tonne-kilometre by 0.32 points for the kite and 0.10 for the rotors, and leave the wing sails' unchanged. Stages \SI{45}{nmi} and points \SI{30}{nmi} apart, at 1.9 times the computing time, raise the cost saving by 0.5 points on average for the kite and 0.25 for the rotors, leave the wing sails' within 0.1 and change the cost of the ship without a rig by at most 0.5\,\%. More headings between the same widest ones, at three times the computing time, lower every ship's cost per tonne-kilometre by about 1.5\,\%, with or without a rig, and move the kite's cost saving by 0.24 points and the others' by less. Each halving of the spacing adds less than the one before, but the kite's gain falls more slowly than the rotors', so even the reference understates the kite more.

\subsection{Forecast studies}
\label{app:forecast_studies}

\tbd{Changing one thing at a time splits the loss of the plan on one forecast run over the 12 departures (\cref{tab:forecast_causes}). On that plan the kite loses $9.6 \pm 2.2$ points of its cost saving, against 5.2 on several members, the wing sails $1.7 \pm 0.3$ and the rotors $5.5 \pm 1.1$, so the kite loses $4.1 \pm 2.5$ points more than the rotors. Split by what the forecast gets wrong, that excess moves between the waves and the wind with the order of the change. Changed first, forecast waves with the real wind cost the kite $4.5 \pm 2.0$ points and take its flying time at sea from 37 to 31\,\%, and cost the wing sails $0.4 \pm 0.3$ and the rotors $-0.1 \pm 0.3$, since only the kite waits on the sea to launch; the wind forecast added next costs the kite no more than the rotors. Changed first instead, forecast wind with the real waves costs the kite $6.9 \pm 1.2$ points, the wing sails $1.8 \pm 0.5$ and the rotors $3.7 \pm 0.6$, an excess over the rotors of $3.2 \pm 1.1$, and the wave forecast added next costs the kite $2.0 \pm 1.9$ and the rotors $1.8 \pm 1.0$. The two errors overlap for the kite: each costs it less once the other is in place, and neither alone carries the excess. Where the wind error sits near the kite's limits cannot be separated from the wind error at large. The kite's lowest flying wind is an apparent wind, which depends on the course and the ship's speed, so no band of true wind picks it out: the 90\,\% of the real \SI{10}{m} winds at samples within \SI{4}{m/s} of that limit lie between 1.8 and \SI{11.1}{m/s}, and that band holds 82\,\% of all samples. After the wave forecast, forecast wind only inside that band costs the kite $3.2 \pm 1.6$ points, the wing sails $1.2 \pm 0.4$ and the rotors $3.6 \pm 0.7$, and the wind outside it adds $1.2 \pm 1.6$ for the kite and $1.9 \pm 0.7$ for the rotors. Split by what the plan decides, the same excess over the rotors is $0.1 \pm 0.2$ points for re-planning, $2.1 \pm 0.9$ for the speed, $-0.1 \pm 0.5$ for the kite's flying decisions and $2.0 \pm 2.4$ for the course when the course changes first, and $0.4 \pm 0.4$ for re-planning, $2.4 \pm 1.5$ for speed and flying together and $1.3 \pm 2.1$ for the course when it changes last. The two splits cut one loss two ways, so the wave row overlaps the speed row: a wave error reaches the cost through the course, the speed and the launches the plan chooses. Whether the kite goes up or stays down, decided on the forecast once course and speed are set, costs $-0.1 \pm 0.5$ points, so the fly-or-stay-down rule is not where the loss lies. At this draft a plan can change only where a stage of the route ends, about every three days in the open ocean, so re-planning once a day is re-planning at that spacing. The speed is the forecast's doing, acted on through a rule of the model: the ship holds the plan's speed, or the next lower one it can hold, when the pull the plan expected does not come. With its course and speed both taken from the forecast the kite flies 26\,\% of the time at sea, against 32\,\% with only the course. A route picked for cost sails the kite close to the wind it needs. On the forecasts $46 \pm 5$\,\% of its time aloft, per departure, lies within \SI{4}{m/s} of the lowest apparent wind it flies in, the forecast's root-mean-square error at four days (\cref{fig:forecast}), and $7 \pm 2$\,\% within the \SI{1.25}{m/s} of the error at six hours. Of all its time at sea it is held down $17 \pm 5$\,\% by the sea, in the earlier model behind these forecast studies, which keeps it down in waves above a fixed height and which the final run replaces with the extra launch wind of \cref{eq:sea_wind}, is held down $25 \pm 3$\,\% by a wind outside its launch range, and stays down $30 \pm 4$\,\% though wind and waves would allow a launch. A deck rig has no such threshold: in less wind it pulls less, but it still pulls. Routing on foresight accounts for part of the kite's loss but not for its excess over the rotors. On one fixed course for all rigs, that of the ship without a rig with perfect foresight, the kite loses $5.8 \pm 1.0$ points against the $9.6 \pm 2.2$ when routed, so routing on the forecasts adds $3.8 \pm 2.1$ points to its loss and $2.8 \pm 1.2$ to the rotors'. There the rotors lose $2.7 \pm 0.5$, the kite still $3.0 \pm 1.2$ more, and the kite keeps $8 \pm 13$\,\% of its cost saving against $78 \pm 5$. Of the kite's routed loss, re-planning once a day costs $0.7 \pm 0.4$ points, the forecast model differing from ERA5 even at the shortest lead $2.5 \pm 1.0$ against $0.9 \pm 0.3$ for the rotors, and the error growing with lead time $6.4 \pm 1.8$ against $4.6 \pm 1.2$. Two things stay open. How a rig's loss divides between course and speed, and between the wave and the wind forecast, depends on which is changed first, by up to 2.5 points of the kite's loss and 1.9 of the rotors', and the wave forecast's part of the kite's excess runs from $4.6 \pm 2.0$ points changed first to $0.2 \pm 2.5$ changed last. And the run cannot separate a wind error near the kite's limits from the wind error at large, nor the pull lost below the wind limit from the time spent down. All these numbers are draft, from a run on e-ammonia.} \tbd{Why the kite flies less and sails longer on forecasts shows in the hours of its ship (\cref{tab:forecast_hours}, paired over the 12 departures). A plan on one forecast run treats that run as certain. The kite then flies $76 \pm 27$ hours less a round trip than with perfect foresight, about what the plan wants up and does not fly, $70 \pm 12$ hours, because the weather met does not allow or pay for it: the effect of a forecast a little off near the limits. Those hours cost little: where the kite does fly it saves as much fuel as before, $-30 \pm 180$~GJ. The cost lies in a passage $80 \pm 39$ hours longer. On the same plans the ship without a rig sails $12 \pm 11$ hours longer and burns $1 \pm 2$\,\% more fuel per tonne-kilometre, while the kite ship sails $65 \pm 36$ hours longer than that and burns $25 \pm 6$\,\% more, so the loss is the kite's own, though the ship with rotors also sails $63 \pm 17$ hours longer on one run. Of the 80 hours, $74 \pm 41$ come from the plan: $134 \pm 110$ nautical miles more distance, about 14~hours, and commanded speeds lower by 0.6~knots. Only $6 \pm 4$ hours are the kite stepping down to a speed it can hold when the planned pull does not come. The plan commands the fast sailing that only the kite can hold for 58~hours less, 80 instead of 138, and holds 44 instead of 117, though point by point the forecast allows it for 290 hours against 308 in the real weather. A commanded speed holds for a whole routing edge, and a flying kite can hold fast sailing over a whole edge far less often: 52 hours on the forecast, against 86 with perfect foresight, $34 \pm 12$ fewer, and the real weather along the course the forecast chose allows no more, 51 hours. The plan commands fast sailing for about 1.5 times those hours on both, so it uses the edges it has as far as the plan with perfect foresight does; a plan on one forecast run finds fewer of them, and the run does not say how much of that is the course it picks and how much the weather the forecast shows. The extra time at sea is spent with the kite down, $156$~hours more, of which $54 \pm 16$ with the wind too far ahead for the kite to pull the ship forward, $46 \pm 38$ through sea and $23 \pm 29$ through too little wind. With perfect foresight it is down 269 hours a round trip for too little wind, 198 of them below the \SI{4.5}{m/s} of apparent wind it needs to stay up, where a following wind cancels the ship's own, and most of the rest between that and the \SI{6}{m/s} from which a pattern pulls; a further 80 hours the wind is enough but from too far ahead. Slowing down raises the apparent wind in a following wind, but a routing speed holds for a whole edge, a day or more, so we also ask of each stretch of these hours, in the weather met, whether one speed held over only that stretch would get the kite up and pulling and whether that pays in fuel, rig wear and the time the slower ship loses, launch included, against the cheapest speed with the kite down there. With perfect foresight a speed exists for 227 of the 269 hours and it pays for 22, 10 of them only as a change within the edge; on one forecast run it pays for 79 of 291, 43 of them only within the edge. Sailed so, these stretches save $0.04 \pm 0.01$\,\% of the fuel, rig wear and time of the round trip as sailed with perfect foresight, $0.30 \pm 0.11$\,\% on one forecast run and $0.15 \pm 0.08$\,\% on several members, so the speed holding for a whole edge costs the kite little. What the kite loses is fuel burned while it is down, $1420 \pm 360$~GJ more a round trip. A plan that keeps a margin inside the limits, in that earlier model with the winds and waves it allows 20\,\% tighter than the real ones (\cref{app:forecast_router}), removes the forced landings and cuts the hours wanted up and not flown to $37 \pm 6$, but the kite then flies $33 \pm 21$ hours less, the ship is no quicker, and in \cref{tab:forecast_plans} the kite keeps $27 \pm 11$\,\% of its cost saving against $26 \pm 10$\,\% on the one run, a gain of $1 \pm 11$ points, so a margin on the limits shows no gain clear of its error. Weighing the members one by one recovers part of the loss: the ship sails $39 \pm 23$ hours less than on the one run, with $27 \pm 24$ hours more fast sailing commanded and the forced landings gone, and the kite keeps $52 \pm 11$\,\%, $26 \pm 12$ points more than on the one run. A plan on the ensemble mean keeps $44 \pm 6$\,\% and one with a margin for wind on the members' quartiles $45 \pm 18$\,\%, gains of $18 \pm 10$ and $20 \pm 16$ points, neither twice its error. The same plan helps the deck rigs about as much: the rotors move from 66 to 79\,\%, a gain of $13 \pm 6$ points, and the wing sails from 72 to 85\,\%, $13 \pm 7$ points. On every plan on forecasts the rotors lead the kite, by 4.9 to 7.5 points of the cost per tonne-kilometre. Holding the kite where its plan put it instead of letting it fly the weather met, 7.8 launches a round trip instead of 7.4, costs $0.7 \pm 0.2$ points of its cost saving. What remains on several members is not knowing: the ship still sails $41 \pm 19$ hours longer than with perfect foresight and the kite flies $56 \pm 28$ hours less, part of which the plan on the finer grid recovers.}

\tbd{Re-planning more often and adapting within the day are measured on their own (\cref{tab:forecast_replanning}, 12 departures), on a route grid eight times as fine as the draft's, where a plan can change at a node about every 10 hours. The wind reforecast starts once a day, which sets its re-planning interval, so this study takes the wind from the operational forecasts issued four times a day, GEFSv11, one model version before the GEFSv12 of the reforecast (\cref{app:forecast_router}); the waves stay on the daily reforecast, driven by GEFSv12, so a fresher wave forecast is not measured and the kite, whose launches wait on the sea, may gain from one more than the table shows. A ship told to re-plan 24, 12 or 6 hours after its last plan can do so only at a node, about every 29, 18 and 10 hours, spacings measured on three departures of the ship without a rig and used here as labels for every rig. On plans about 29 hours apart the operational forecasts and the reforecast agree within the error, $20 \pm 8$ against $15 \pm 9$\,\% of the kite's cost saving, $75 \pm 6$ against $70 \pm 6$ of the wing sails' and $84 \pm 4$ against $84 \pm 3$ of the rotors', a check that spans the change of model version and pairs GEFSv11 wind with waves from GEFSv12, so it shows that these two archives agree here, not that the version does not matter. Re-planning more often on the operational forecasts helps the wing sails most: about every 10 hours they keep $85 \pm 4$\,\%, $10 \pm 4$ points more than about every 29 hours. It is the largest of the 18 paired changes of \cref{tab:forecast_replanning} against its error; among 18, one this far from zero arises by chance about one time in six if the errors are normal, and more often on 12 departures, so we read it as a likely gain rather than a proven one. The interval and the forecast's age change together, since no case re-plans about every 10 hours on a forecast issued once a day, so the table does not say how much of the gain is the forecast being fresher. The kite gains $6 \pm 7$ points and the rotors $1 \pm 2$, neither clear of its error, and at about 18 hours the kite is $3 \pm 4$ points lower, inside it. A ship that also adapts between plans chooses each step's speed, or also its next node, in the forecast corrected by the error it meets at its own position, held over the step, plus the router's estimate of the cost of the rest of the passage from where the step ends. Chosen so, the ship without a rig pays 0.1 to 0.3\,\% less per tonne-kilometre, within the error of $0.3$, and no rig gains clear of its error: on plans about every 29 hours, adapting the speed keeps the kite $10 \pm 8$ points more, the course too $8 \pm 10$, and costs the wing sails $5 \pm 6$. Against plans about every 10 hours without adapting, adapting the speed and the course costs the wing sails $15 \pm 7$ points, about twice its error, the kite $3 \pm 10$ and the rotors $1 \pm 4$. That estimate is a coarser search of the corridor than the route itself and can lie on either side of the true cost, and how much of these losses comes from its error is not measured. The same plan on the finer grid, the one forecast run re-planned about every 29 hours, keeps $15 \pm 9$, $70 \pm 6$ and $84 \pm 3$\,\% against $-8 \pm 20$, $74 \pm 5$ and $65 \pm 6$\,\% on the draft's grid in the model of that study, which folds a finer grid and more frequent re-planning together, so the final run on a finer grid will move the shares above. Re-planning once a day, which is what the reforecast allows, may therefore understate what the wing sails keep by up to 10 points at these departures, and leaves what the kite and the rotors keep unchanged within the error; the study supports no claim that adapting within the day recovers more than the plan.}

\begin{table}[tbp]
\centering
\caption{Share of the cost saving with perfect foresight on the draft's route grid that each rig keeps when re-planned every 24 hours on each plan, on the archived forecasts or, where the row says so, on the real weather, against the ship without a rig on the same plan, for the \SI{140}{m} e-ammonia ship to Guadeloupe, with its standard error over the departures [\%]. \tbd{Draft, 12 departures.} Method study, \studyrun{forecast_plans}. \stale{forecast_plans}{}}
\label{tab:forecast_plans}
\InputIfFileExists{sections/forecast_plans_table}{}{\tbd{Forecast plans table pending: run the forecast\_plans\_table step.}}
\end{table}
\begin{table}[tbp]
\centering
\footnotesize
\caption{Share of the cost saving with perfect foresight each rig keeps per re-planning case, for the \SI{140}{m} e-ammonia ship to Guadeloupe on a route grid of \SI{90}{nmi} stages, against the ship without a rig on the same plan, with its standard error over the departures [\%], and the change from the plan the row names or, where it names none, from the operational forecasts about every 29 hours, paired per departure [points]. A plan changes at the first node 24, 12 or 6 hours after the last plan, about every 29, 18 and 10 hours as measured on three departures of the ship without a rig. \tbd{Draft, 12 departures.} Method study, \studyrun{forecast_replanning}.}
\label{tab:forecast_replanning}
\InputIfFileExists{sections/forecast_replanning_table}{}{\tbd{Forecast re-planning table pending: run the forecast\_replanning\_table step.}}
\end{table}
\begin{table}[tbp]
\centering
\footnotesize
\caption{Hours and fuel of the kite ship per round trip with perfect foresight, re-planned every 24 hours on one forecast run, on several forecast members, and on one forecast run in a plan with the kite's wind limits 20\,\% tighter and its extra launch wind in a sea 20\,\% higher than the real ones, for the \SI{140}{m} e-ammonia ship to Guadeloupe, mean with standard error over the departures. Fast sailing is at 12 knots or more, which only a flying kite holds on this ship; ``possible'' says where the real weather let it hold and ``forecast'' where the forecast the ship planned on said it could, sample by sample and over edges where every sample allows it. On several members, what the plan wanted and what the forecast said are read on the first member. \tbd{Draft, 12 departures.} Method study, \studyrun{forecast_hours}.}
\label{tab:forecast_hours}
\InputIfFileExists{sections/forecast_hours_table}{}{\tbd{Forecast hours table pending: run the forecast\_hours\_table step.}}
\end{table}
\begin{table*}[tbp]
\centering
\footnotesize
\caption{Where each rig's cost saving goes when it is re-planned every 24 hours on one run of the archived forecasts, for the \SI{140}{m} e-ammonia ship to Guadeloupe, and the kite's excess over the wing sails and over the rotors, paired per departure, each with its standard error over the departures. Cost savings are against the ship without a rig on the same plan. Each block splits the loss by changing one thing at a time: the plan's course, then re-planning once a day, then its speed, then the kite's decision to fly (course first); re-planning, then speed and kite together, then the course (course last); re-planning, the forecast model against ERA5 at its shortest lead, then the error that grows with lead time; and, after re-planning, the wave and the wind forecast in both orders, the wind split with the waves first by whether the wind speed lies in the band near the kite's lowest flying wind. The last block is the kite's time near its limits on the forecasts per departure. The fixed course is the one the ship without a rig takes with perfect foresight. \tbd{Draft, 12 departures, in the earlier model with a fixed wave limit for a launch.} Method study, \studyrun{forecast_causes}.}
\label{tab:forecast_causes}
\InputIfFileExists{sections/forecast_causes_table}{}{\tbd{Forecast causes table pending: run the forecast\_causes\_table step.}}
\end{table*}

\subsection{Model checks}
\label{sec:checks}

These checks say what the model reproduces and where it departs from measurements.

\begin{itemize}
\item \textbf{Force chain worked by hand.} \Cref{tab:hand} follows the wind through shear and the ship's speed to the apparent wind at the rig and the force along the course, calculated without the code. Thrust, side force and spin power agree to the digits printed. The massless \SI{800}{m^2} kite in \SI{12}{m/s} of apparent wind pulls \SI{64.0}{kN} by hand and in the code, and its thrust falls to nothing at an apparent wind angle of \SI{15}{\degree}, the upwind limit of its glide ratio. Taking the kite's wind at an elevation of 15 or \SI{8}{\degree} instead of \SI{25}{\degree} lowers its cost saving by \tbd{0.4 to 2.9} points on the \SI{140}{m} ship, and half rather than a quarter of the tether's drag acting at the kite changes it by \tbd{$-2.0$ to $+0.6$} points.
\item \textbf{Calm-water resistance.} The wetted surface, friction coefficient and wave-resistance coefficients of the worked example of \citet{holtrop1982} come out within 0.2\,\%, its bulb resistance within 5\,\% and its correlation allowance within 0.5\,\%.
\item \textbf{Calibrated class power.} Every standard class except the coaster needs the calm-water delivered power of its source design, through a resistance factor $f_R$ of 0.61 to 0.99 (\cref{app:calm_calibration}). Uncalibrated, on the worked \num{50000}~t MR tanker of \citet{man2011}, the model needs 7\,\% more power than MAN's light propeller curve at \SI{14}{kn}.
\item \textbf{Coaster power.} The \SI{85}{m} coaster gets \SIrange{700}{1500}{kW} installed. The Damen Combi Freighter 3850, \SI{89.7}{m} overall and carrying \num{3830}~t, has a \SI{1104}{kW} main engine for \SI{11.8}{kn} \citep{maritimeexecutive2022}.
\item \textbf{Added resistance.} The non-dimensional added resistance peaks between 5 and 10 in head waves and stays within $\pm2$ in following waves, inside the scatter of the model tests behind \citet{mittendorf2022}.
\item \textbf{Kite without mass.} The quasi-steady kite without mass carries the tension of Eq.~39 of \citet{loyd1980} with zero weight, at the roll angle of his Eqs.~40--41, to a relative error below $10^{-6}$.
\item \textbf{Fuel saving against operation.} A kite of \SI{125}{m^2} projected, \SI{160}{m^2} flat, on a \SI{130}{m} VLSFO ship, the size the \SI{132}{m} Beluga SkySails was first launched with \citep{wiwo2008,noaa2009}, saves \tbd{1.7}\,\% of fuel on the routed round trip to Guadeloupe\tbd{, flying 6\,\% of the time at a cost saving of $-0.3$\,\%. Read as \SI{160}{m^2} projected it saves 3.2\,\%}. The WINTECC record reports fuel savings of 5\,\% on the average route mix and 10 to 12\,\% on North Atlantic and North Pacific routes for that ship \citep{wintecc2009_record}, without giving the kite size. \tbd{Removing the flight drag gives 1.6\,\%, a force limit of \SI{0.9}{kN/m^2} 1.5\,\%, and free canopy wear and launches 2.2\,\%, all within the routing noise.} We routed one route, not the reported mix, and did not find what explains the gap. The model comes out below the report although its router sees the weather ahead, so this check does not bound our savings from above.
\item \textbf{Routing noise.} A change of 2\,\% in the tether length or the smallest turn radius should barely move a kite's cost saving. \tbd{On three ships with kites of \SI{3200}{m^2}, \SI{1448}{m^2} and \SI{16000}{m^2}, the base and those four changes span 1.4, 0.8 and 1.5 points of cost saving and 6.6, 4.9 and 2.3 points of fuel saving, in the method study \studyrun{base_kite}.} The router minimises cost, so routes of nearly equal cost can burn different amounts of fuel. We do not read a smaller difference between two kites as a result.
\item \textbf{Routing grid.} Doubling the spacing of stages and points lowers the cost saving per tonne-kilometre by 1.0 points on average for the kite, 0.7 for the rotors and 0.5 for the wing sails, on the 12 departures of two ships on three routes. \Cref{sec:routing} gives the effect of halving it.
\end{itemize}

\subsubsection{Error of the missed-goal figures}
\label{sec:check_missed}

All limits come from one size sweep over the same \tbd{6} departures, around the size the class study gives each rig: half to twice for wing sails and rotors as far as they fit on deck, and a quarter to four times for the kite up to the search bound of \SI{16000}{m^2}. \tbd{The best kite lies inside the sizes swept on all three ships. The best wing sails are the largest that fit on deck on the \SI{140}{m} ship and the Neo-Panamax and the best rotors on the Neo-Panamax. The rest lie inside the sizes swept.}

Routing again at each canopy life moves the limits little: on the Neo-Panamax on e-ammonia to New York the limit with the route held is \tbd{\num{1610}} flight hours and the one routed afresh at each life \tbd{\num{1600}}.

\Cref{sec:res_missed} gives what staying at today's canopy price takes from the kite on three ships. On three departures of a draft, the routes changing between the two prices moved these figures by 0.1 points at most. A line through the prices from USD~112 to 375/m$^2$ gave \tbd{1.6, 1.6 and 0.7} points against \tbd{1.7, 1.6 and 0.8} between the two prices. Leaving out of maintenance the canopy and tether that wear already pays for was worth \tbd{0.3, 0.15 and 0.24} points of cost saving on the three ships at today's price.

\begin{table}[tbp]
\centering
\caption{The chain from true wind to force on the ship, worked by hand independently of the code and from the code, at a ship speed of \SI{5.4}{m/s} and a wind shear exponent of 0.05. Wing sails of \SI{800}{m^2} at $C_L$ 1 and rotors of \SI{400}{m^2} at spin ratio 3 in a wind of \SI{10}{m/s} at \SI{10}{m}. The kite of \SI{800}{m^2}, at its reference height of \SI{269}{m} and with no side weight, in winds at \SI{10}{m} of 5.88, 6.62, 8.83, 11.16 and \SI{13.80}{m/s}, which put its apparent wind on nodes of its pattern table; the hand calculation takes the force coefficients there from the table. True wind angle TWA, apparent wind $V_a$ and its angle AWA, thrust $F_x$ and side force $F_y$.}
\label{tab:hand}
\adjustbox{max width=\columnwidth}{%
\begin{tabular}{@{}lrrrrrrr@{}}
\toprule
Rig & TWA & $V_a$ & AWA & \multicolumn{2}{c}{$F_x$ [kN]} & \multicolumn{2}{c}{$F_y$ [kN]} \\
\cmidrule(lr){5-6}\cmidrule(lr){7-8}
 & [\si{\degree}] & [m/s] & [\si{\degree}] & hand & code & hand & code \\
\midrule
Wing & 90 & 11.69 & 62.5 & 51.76 & 51.76 & 41.42 & 41.42 \\
Wing & 45 & 14.69 & 29.9 & 35.14 & 35.14 & 100.99 & 100.99 \\
Wing & 150 & 6.30 & 124.6 & 17.23 & 17.23 & $-8.24$ & $-8.24$ \\
Wing & 180 & 4.97 & 180.0 & 2.05 & 2.05 & $-12.11$ & $-12.11$ \\
Rotor & 90 & 11.76 & 62.7 & 243.77 & 243.77 & 227.31 & 227.31 \\
Rotor & 45 & 14.77 & 30.0 & 131.59 & 131.59 & 510.60 & 510.60 \\
Rotor & 150 & 6.37 & 124.9 & 91.93 & 91.93 & $-32.69$ & $-32.69$ \\
\midrule
Kite & 26.6 & 12.00 & 15.0 & $-0.03$ & $-0.03$ & 63.73 & 63.73 \\
Kite & 50.2 & 12.00 & 30.0 & 19.48 & 19.48 & 106.00 & 106.00 \\
Kite & 86.7 & 12.00 & 60.0 & 113.45 & 113.45 & 249.73 & 249.73 \\
Kite & 114.2 & 12.00 & 90.0 & 241.07 & 241.07 & 216.84 & 216.84 \\
Kite & 148.6 & 12.00 & 135.0 & 337.74 & 337.74 & 77.98 & 77.98 \\
\bottomrule
\end{tabular}}
\end{table}

\subsection{Model detail against its effect}
\label{sec:model_detail}

\Cref{tab:model_detail} lists every part of the model with how it is modelled, how much it moves each rig's cost saving and which rig is cheapest, and its known error. A part the model can do without is switched off on the three ships of the development study (\cref{sec:res_inputs}): the calibration of the calm-water resistance, the added resistance in waves, the windage, the heel, leeway and rudder limits, and the rudder's drag. A part that development can change is moved over the ranges of its inputs in that study, one at a time; the propeller and the wind at height are not, so their effect is not measured here. \tbd{All at draft fidelity, where the routing noise of a kite's cost saving is up to 1.5 points (\cref{sec:checks}).}

\tbd{No part is detailed far beyond its effect. The heel, leeway and rudder limits move no cost saving by more than 0.4 points, and we keep them because they cost nothing and stop the router from using trims a ship could not hold. Two rounds of the size search lose at most 0.2 points against every size. The rudder drag costs rotors in a row up to 1.3 points. Towing the kite from the bow mast rather than the hull's lateral centre moves its cost saving by 0.1 points on average (\cref{sec:tow_point}). Each weather year alone moves a cost saving by up to 1.7 points, so a longer record would narrow the results without changing which rig is cheapest. Two parts change the cheapest rig or move a saving by several points and are coarser than their effect warrants, each a candidate for a better model: the extra wind a kite's launch needs in a sea, which moves its cost saving by up to 3.6 points and rests on one reported pair of wind speeds (a launch envelope measured at sea), and the corridor spacing (a finer or adaptive corridor), which moves the kite's cost saving by up to 1.8 points and decides between the kite and the rotors on the \SI{140}{m} ship, where they lie within 0.7 points of each other. On the corridor of the results the kite's saving still lies 1.3 points and the rotors' 0.5 below a reference finer in both, so the router is not converged and its error favours the rotors (\cref{sec:routing}). The one mean power-law wind profile at the kite's height is coarse as well, and its effect is not measured here (better: hourly profiles from the ERA5 model levels).}

{\footnotesize
\setlength{\LTcapwidth}{\textwidth}
\setlength{\tabcolsep}{3.5pt}
\setlength{\extrarowheight}{1pt}
\begin{longtable}{@{}>{\raggedright\arraybackslash}p{2.5cm}>{\raggedright\arraybackslash}p{3.9cm}rrr>{\raggedright\arraybackslash}p{2.1cm}>{\raggedright\arraybackslash}p{3.5cm}@{}}
\caption{Each part of the model, how its effect is measured, its model, its effect on the cost saving of each rig, the ships where it changes which rig is cheapest, and its known error. The effect is the largest difference in a rig's cost saving, in percentage points, between the model as it stands and the model with the part changed, or with the inputs named at the part moved over their ranges, over the three ships of the development study. The parts switched off and the weather years come from the method study \studyrun{model_parts}, the router from \studyrun{convergence} and the size search from \studyrun{size_search}.}
\label{tab:model_detail}\\
\toprule
Part & Model & \multicolumn{3}{c}{Saving moved [pt]} & Cheapest rig changes & Known error \\
\cmidrule(lr){3-5}
 & & kite & wings & rotors & & \\
\midrule
\endfirsthead
\multicolumn{7}{@{}l}{\cref{tab:model_detail} (continued)}\\
\toprule
Part & Model & kite & wings & rotors & Cheapest rig changes & Known error \\
\midrule
\endhead
\bottomrule
\endfoot
\InputIfFileExists{sections/fidelity_table}{}{\multicolumn{7}{@{}l}{\tbd{Pending: run the fidelity\_table step.}}\\}
\end{longtable}}

\subsection{Where each rig acts on the ship}
\label{sec:tow_point}

A rig's side force turns the ship unless it acts where the hull's own side force acts, at the hull's lateral centre about a third of the length forward of midship, and the rudder holds any moment left over at the cost of drag. A single wing sail or rotor stands at that centre, while a row is centred at midship. The literature differs on how much this matters: on a cargo ship modelled from CFD the position of the sails changed the resistance by only 1 to 2\,\% \citep{kramer2022}, while ShipCLEAN gains several points of fuel saving by moving rotors aft \citep{tillig2020}. We use the effective rudder aspect ratio of about 1.6 that \citet{kramer2022} derive from CFD, rather than the doubling often assumed because the hull above the rudder acts as a mirror plane.

The kite's side force acts where its tether meets the ship. Ahead of the hull's lateral centre it turns the bow to leeward, aft of it the stern, and the rudder holds the difference at the cost of its drag (\cref{sec:balance}). We tow every soft kite from the foot of its launch mast at the forward end of the deck, ahead of the lateral centre by 2 to 10\,\% of the length. To see what that costs, we flew each class's best kite on e-ammonia with its tow point at the bow, the lateral centre and midship, on the westerly and trade-wind routes (\cref{fig:tow_point}). Only the rudder load changes between them in the model.

\begin{figure*}[tbp]
\centering
\plot[\textwidth]{tow_point}
\caption{Fuel (a) and cost (b) saving of each class's kite with its tow point at the hull's lateral centre, at midship with the mast taking the deck slots it takes at the bow, the least it can take there, and at midship with the mast taking one whole deck bay of a container ship, the most, less the same saving with its tow point at the bow, the line at zero; the marker is the mean over the two routes and the bar their range. Classes are sorted by deadweight. Method study, \studyrun{tow_point}.}
\label{fig:tow_point}
\end{figure*}

The bow costs almost no thrust. The lateral centre saves \tbd{0.02} points more fuel than the bow on average, less than the spread between ships and routes as the router takes a different path, and midship \tbd{0.7} points less, as its side force acts about a third of the length aft of the lateral centre. The cost saving moves by \tbd{0.1} points on average and at most \tbd{1.5}.

Cargo space decides the rest. On a bulk carrier or tanker a mast at midship takes no cargo space but gains nothing in thrust either. On a container ship it stands in the deck stack, and since no container ship flies a kite we estimate the loss from the geometry: at most one 40-foot bay across the beam on every tier on deck, and at least the single row of a bay that the bow mast takes (\cref{sec:model}). With the bay and the rudder load together, midship saves \tbd{1.3 to 3.5} points less cost than the bow on the container ships that fill their slots before their deadweight. We therefore keep the bow, as the kites built so far do \citep{fritz2013}, and give up the \tbd{0.02} points of fuel saving of the lateral centre, which lies within the cargo area on most ships.

\subsection{Kite flight and launch}
\label{app:kite_flight}

\Cref{tab:flight} gives where the kite works by band of wind angle (\cref{sec:res_forces}), and \cref{tab:sea_effect} its cost saving with and without the extra launch wind in a sea (\cref{sec:kite,sec:res_inputs}).

\begin{table*}[tbp]
\centering
\caption{Where the kite works over the \tbd{28} cases with a kite, each with its cheapest size, by band of true wind angle: share of the time in the band, share of that time in which the kite could pull, is up and pulls, its mean thrust while up, lulls and upwind drag included, the share of the hours at each ship speed in which it could pull and in which it is up, the speed of the ship without a rig and of the kite ship in that band, and the lowest apparent wind angle (AWA) at which each rig pulls. \tbd{Draft run.}}
\label{tab:flight}
\adjustbox{max width=\textwidth}{\InputIfFileExists{sections/flight_table}{}{\tbd{Flight table pending: run the flight\_table step.}}}
\end{table*}

\begin{table}[tbp]
\centering
\caption{The kite's cost saving per tonne-kilometre against the same ship without a rig and its share of the time at sea in the air, with the extra launch wind of \cref{eq:sea_wind} and with no sea effect on the launch, on e-ammonia, averaged over the six routes. \tbd{Draft run; the \SI{140}{m} ship over four routes, as the router found no passage on New York and Houston.} Method study, \studyrun{sea_effect}.}
\label{tab:sea_effect}
\InputIfFileExists{sections/sea_effect_table}{}{\tbd{Sea effect table pending: run the sea\_effect\_table step.}}
\end{table}

\subsection{Interference between deck rigs}
\label{app:interference}

In two-dimensional potential flow, three wing sails in line lose up to 6\,\% of thrust and four in a two-by-two layout 10 to 28\,\% \citep{vanreen2026}. In the wind tunnel, two rotors three diameters apart in line lost up to 37\,\% of driving force at \SI{30}{\degree} apparent wind, gained up to 17\,\% downwind, and lost 4 to 6\,\% on average over the wind angles of a reaching ship \citep{bordogna2020}. CFD of wing sails on a bulk carrier gives 98\,\% of the isolated thrust for three sails 2.4 chords apart and 77\,\% for five sails 1.4 chords apart \citep{mboumba2022}, the chord being the depth of a wing from front to back. Our array factor gives 91 and 85\,\%, which takes back a third of the loss for sails trimmed one by one, and 87 and 78\,\% at fixed trim.

\subsection{Wind shear by region}
\label{app:shear}

In the well-mixed trade-wind belt our ERA5 profiles give only 1.17 times the \SI{10}{m} wind at \SI{540}{m}. North of \SI{40}{\degree N} they follow the logarithmic law with \SI{0.2}{mm} roughness, an exponent of about 0.08 between \SI{10}{m} and \SI{270}{m}, which matches the four-year mean of mast and lidar measurements \SI{85}{km} off the Dutch coast up to \SI{315}{m} better than a power law with the usual offshore exponent of 0.14 \citep{kalverla2020thesis}.

\subsection{Calm-water resistance of full hulls}
\label{app:full_hulls}

\citet{holtrop1982} limit their method to hulls resembling their average ship and call its high-block adjustments tentative. \citet{ortigosa2009} tabulate a prismatic coefficient of at most 0.85 for tankers and bulk carriers, a limit we could read only second-hand. Beyond it the fitted wave term, the part of the resistance due to the waves the hull makes, grows steeply. Uncalibrated it takes 25 and 31\,\% of the resistance of the Ultramax and Kamsarmax at service speed, and one calibration factor scales every component, which is what would make their power at operating speed too low.

\subsection{Cost split against published studies}
\label{app:cost_split}

Published cost splits leave out the interest on the cargo, which the cargo owner bears, and mostly the port dues. On that basis fuel is \tbd{30}\,\% rather than a quarter of the cost of our \SI{85}{m} ship and \tbd{41}\,\% of our \SI{200}{m} ship on VLSFO, at the speeds chosen by cost of about \tbd{10 to 11}~kn. \Citet{karvonen2020} give 23\,\% for a general cargo ship of \num{3500}~t deadweight at \SI{9}{kn} and 51\,\% for a bulk carrier of \num{60000}~t at \SI{11.6}{kn}, with fuel about a sixth cheaper than ours. The day rates of \citet{healy2019} give 30 to 45\,\% for a Handysize bulk carrier between 10 and \SI{13}{kn}, once we count their earnings as the return on capital. Fuel's share grows with speed, so ours sits where these put it. On the same basis capital with maintenance and insurance is \tbd{48 to 50}\,\% of our cost against 34 to 58\,\% in \citet{karvonen2020}, crew and other running costs are \tbd{9 to 17}\,\% against their 11 to 13\,\%, larger on the smaller ship, and port dues are about \tbd{4}\,\%. Cargo interest grows from \tbd{an eighth} of our cost on the \SI{85}{m} ship to \tbd{two fifths} on the \SI{200}{m} ship, because every ship of the design study carries general cargo worth USD~\num{2500} a tonne (\cref{sec:inputs}); a bulk carrier's cargo is worth a tenth of that or less, and the class study prices each class's own cargo (\cref{tab:classes}).

\subsection{Published statements on kites and ship speed}
\label{app:speed_literature}

\citet{cedelft2016} find, at 100 and 85\,\% of service speed, that kites save relatively more at the lower speed and the same or less in absolute terms at the higher, and conclude that ``towing kites are best-suited for lower ship speeds'' (p.~48). \citet{traut2014} assume a 20\,\% speed reduction and \citet{bentin2018} exclude ships of \SI{20}{kn} or more. \citet{emsa2023} state that kites do better at higher speeds, on one provider-reported pair of figures (p.~30).

The kites of \citet{podeur2016} and \citet{leloup2013} pull with a glide ratio of about 4.6 to 4.7, against our 3.9. Podeur's ship moves only along its heading and \citet{leloup2013} take the kite's drift force as negligible, so in both the side force costs nothing, while ours pays leeway, rudder drag and heel for it. The 26.6\,\% of \citet{leloup2013} is one eastbound crossing in one recorded set of winds, and Podeur's eastbound leg saves 18.2\,\% against 6.0\,\% westbound. The routed figures of \citet{bentin2018}, 19 to 29\,\%, let the route grow by up to a fifth and take the kite's pull from an empirical law in the apparent wind.
